\section{Results}

The fitted parameters and errors are summarized in \cref{tab:fit-results}.
The pressure-dependent tangent modulus captures the rapid low-pressure closure
of compliant cracks and the high-pressure plateau of each specimen.  The
serpentinized specimens retain smaller high-pressure drained moduli than HT14,
consistent with mineral softening and reaction-generated microstructure.

\begin{table}[ht]
 \centering
 \caption{Crack-closure parameters and comparison with the released Figure~3c data.}
 \label{tab:fit-results}
 \begin{tabular}{lrrrrr}
  \toprule
  Sample & \(K_\infty\) [GPa] & \(C_c\) [GPa\(^{-1}\)] & \(P_c\) [MPa] & RMSE in \(B\) & max. \(|\Delta B|\) \\
  \midrule
  HT14 & 51.053 & 0.01498 & 13.369 & 0.00532 & 0.00818 \\
  SP14 & 35.994 & 0.12772 & 7.347 & 0.00872 & 0.01962 \\
  SP30 & 24.068 & 0.07160 & 8.760 & 0.00589 & 0.01122 \\
  \bottomrule
 \end{tabular}
\end{table}

\begin{figure}[ht]
 \centering
 \includegraphics[width=0.92\linewidth]{moose_lawal_kim_biot_replication.pdf}
 \caption{Pressure-dependent Biot coefficients calculated in the released
 Lawal--Kim Figure~3c workbook and the constrained implicit-AD constitutive-path
 calculation. The crack-compliance parameters are calibrated to the drained
 tangent moduli from which the workbook coefficients were calculated. The
 curves therefore show calibration and differentiation consistency rather than
 independent prediction.}
 \label{fig:lawal-kim-calibration}
\end{figure}

The MOOSE tangent agrees with the analytic derivative in
\cref{eq:pressure-path-tangent} to maximum absolute errors of
\(1.3\times10^{-9}\), \(1.5\times10^{-8}\), and \(9.8\times10^{-9}\) for
HT14, SP14, and SP30, respectively.  The largest local constraint norm is
\(2.3\times10^{-9}\). These comparisons verify the implicit derivative
producing \(B\) for the calibrated constitutive path. They do not constitute an
independent validation of the pressure dependence.

Over the experimental range, the finite-deformation correction relative to
\(1-K(P)/K_s'\) remains small: its maximum magnitude is approximately
\(2.9\times10^{-4}\) for HT14, \(7.8\times10^{-4}\) for SP14, and
\(7.1\times10^{-4}\) for SP30.  The observed pressure dependence is therefore
controlled primarily by crack closure through \(K(P)\), while the nonlinear
specific-volume factor provides a systematic finite-deformation correction.
